\section{Implementation}\label{sec:implementation}

\sys{} is implemented in approximately 15{,}000 lines of core
Python~3.9 (excluding tests, scripts, and evaluation tooling),
with ARM firmware written in bare-metal C ($\sim$1{,}600 lines across
six device variants).  We highlight key implementation
decisions below.

\subsection{3D Simulation}

We build on Habitat-Lab v0.3.0 with the Habitat-Sim renderer.
Scenes are loaded from the HSSD-Hab and Replica-CAD datasets.
Humanoid navigation uses the built-in shortest-path follower with a
custom task scheduler that dispatches activities from the LLM-generated
schedule.  Object interactions (pick, place, open, toggle) use
Habitat's rearrangement primitives.

\subsection{Firmware Containers}

Each device type has a corresponding Dockerfile that:
\begin{enumerate}
  \item Installs the QEMU ARM system emulator (v10.2).
  \item Copies the pre-compiled firmware ELF binary for the
    target device type.
  \item Exposes a TCP socket (default port~15\,000) for UART I/O.
  \item Launches QEMU with the LM3S6965EVB machine profile:\\
    \texttt{qemu-system-arm -machine lm3s6965evb}\\
    \texttt{-cpu cortex-m3 -nographic}\\
    \texttt{-kernel <device>.elf}\\
    \texttt{-serial tcp::15000,server,nowait}
\end{enumerate}
A Python \texttt{DeviceFirmwareManager} class wraps the Docker SDK to
create, start, stop, and remove containers.  Health checks ping the
TCP socket; unresponsive containers are restarted automatically.
A \texttt{DeviceType} enum with 28~name aliases (e.g., ``thermostat''
$\to$ \textsc{Temperature\_Sensor}, ``bulb'' $\to$
\textsc{Smart\_Light}) enables fuzzy device-type resolution from
natural-language scene descriptions.

\paragraph{Firmware Build System.}
All six firmware variants are compiled with a shared Makefile using
\texttt{arm-none-eabi-gcc} and the flags
\texttt{-mcpu=cortex-m3 -mthumb -Os -nostartfiles -nostdlib
-ffunction-sections -fdata-sections}.  A custom linker script
(\texttt{linker.ld}) maps the ISR vector table to flash address
\texttt{0x0000\_0000} and allocates 64\,KB flash and 20\,KB SRAM.
Each firmware image is approximately 300~lines of C implementing
the UART driver (PL011-compatible registers at \texttt{0x4000C000}),
command dispatcher, GPIO handlers, and device-specific sensor models.
The build produces both ELF (for QEMU) and raw binary outputs.

We provide reference firmware images for six device types:
\begin{itemize}[leftmargin=*]
  \item \textbf{Smart Light} --- PWM-controlled LED with on/off/dim,
    brightness (0--100\%), and color temperature
    (2\,700--6\,500\,K).
  \item \textbf{Motion Sensor} --- PIR interrupt with configurable
    sensitivity (1--10) and cooldown timer.
  \item \textbf{Temperature Sensor} --- ADC sampling at 1\,Hz with
    Kalman filtering and threshold alerts.
  \item \textbf{Humidity Sensor} --- combined RH\% and temperature
    monitoring with user-invocable calibration.
  \item \textbf{Door Sensor} --- magnetic contact with open/closed
    state and tamper detection.
  \item \textbf{Smart Plug} --- relay-driven outlet with real-time
    power metering (watts) and schedule control.
\end{itemize}

\paragraph{UART Communication Protocol.}
Each firmware processes newline-terminated ASCII commands over UART0.
The host-side Python driver sends commands via a TCP socket and
accumulates byte-by-byte QEMU output (a consequence of the
hardware-accurate UART emulation) into complete response lines.
The shared command set includes: \texttt{IDENTIFY}, \texttt{ON/OFF},
\texttt{STATUS}, \texttt{GET\_ALL}, \texttt{AUTH:<token>},
\texttt{SET\_ID:<name>}, \texttt{FW\_UPDATE:<hex>}, \texttt{FW\_APPLY},
\texttt{DEBUG\_DUMP}, \texttt{REBOOT}, \texttt{SET\_SCHEDULE}, and
\texttt{CALIBRATE}, plus device-specific extensions (e.g.,
\texttt{SET\_BRIGHTNESS}, \texttt{SET\_SENSITIVITY}).

\subsection{LLM Integration}

\sys{} communicates with LLMs through the OpenAI-compatible API
exposed by LMStudio on \texttt{localhost:1234}.  The
\texttt{LLMClient} class supports configurable model selection,
temperature, and retry logic.  Prompts are Jinja2 templates
parameterized by persona fields, day type, and room inventory.

All experiments and the large-scale autonomous evaluation use
GPT-OSS~20B, a 20-billion-parameter open-source model running
locally with 4-bit quantization (GGUF format) on consumer hardware
(Apple M2~Pro, 32\,GB RAM).  The model generates structured JSON
schedules conditioned on persona profiles, room inventories, and
day type (weekday vs.\ weekend).

\subsection{SmartThings Integration}

The Schema Connector is a Flask application that handles the
SmartThings lifecycle webhooks.  OAuth2 tokens are managed via a
token store backed by SQLite.  Device profiles are auto-generated
from the Device Registry, mapping \sys{} device types to
SmartThings capability identifiers (e.g.,
\texttt{st.capability.switch}, \texttt{st.capability.temperatureMeasurement}).

State updates are pushed to SmartThings via HTTPS POST to the
state callback URL, using exponential backoff on transient failures.
Incoming commands are translated to firmware-level GPIO events and
dispatched through the event bus.

\subsection{Simulated Home Network}

The simulated network is implemented as a configurable Docker network
(\texttt{vesper-home-network}, \texttt{172.20.0.0/24}) supporting
four driver modes via a \texttt{NetworkMode} enum:
\texttt{BRIDGE} (default), \texttt{MACVLAN}, \texttt{IPVLAN}, and
\texttt{HOST}.  A \texttt{NetworkConfig} dataclass centralizes all
topology parameters: gateway (\texttt{172.20.0.1}), MQTT broker
(\texttt{172.20.0.2}), device IPs (starting at \texttt{172.20.0.10}),
latency (5.0\,ms $\pm$ 2.0\,ms jitter), bandwidth (100\,Mbps), and
packet-loss ratio.

\paragraph{Wireshark Integration.}
Attack traffic is captured through two complementary paths.
First, \texttt{tshark}~4.6.3 (Wireshark CLI) runs on the loopback
interface with a BPF filter targeting the firmware TCP ports,
writing per-attack and global \texttt{.pcap} files that are
independently verifiable with \texttt{wireshark} or
\texttt{tcpdump} (\Cref{sec:eval-pcap}).
Second, an \texttt{InstrumentedSocket} wrapper records every
TX/RX payload with nanosecond timestamps at the application layer,
providing structured packet logs even when Wireshark is unavailable.
When macvlan or ipvlan mode is selected, a
\texttt{WiresharkLiveCapture} module can additionally start
Wireshark on the parent NIC (auto-detected via \texttt{ifconfig}
on macOS or \texttt{ip route} on Linux), enabling real-time
inspection of raw container frames.
This multi-path architecture ensures that attack traffic evidence
is always available (\Cref{tab:traffic-analysis}).

\paragraph{MQTT Broker.}
The \texttt{MQTTBrokerSimulator} is a threaded Python server
accepting up to 32~concurrent client connections.  It implements
publish, subscribe (including wildcard \texttt{\#} and \texttt{+}
matching), retained messages, and optional username/password
authentication.  The broker routes device telemetry through a
hierarchical topic scheme:
\texttt{vesper/devices/<id>/\{state,events,commands\}}.

\paragraph{Protocol Simulators.}
Pluggable \texttt{ProtocolSimulator} instances provide lightweight
application-layer models of Zigbee, Z-Wave, and BLE communication.
Each simulator handles device registration, message routing with
configurable latency and packet loss, and network-key management.
Protocol-specific RF framing is not modeled; instead, simulated
protocol attacks (Zigbee replay, key extraction) operate over TCP
connections to the firmware containers.  A \texttt{PacketCapture}
module records all traffic with nanosecond timestamps for post-hoc
forensic analysis.

\subsection{Attack Framework}

The security testing pipeline consists of five complementary attack
suites totaling 36~unique attacks in approximately 3{,}200~lines
of Python.

\paragraph{Firmware Attacks (Suite~1).}
The \texttt{FirmwareAttackFramework} class implements 18~firmware-level
attacks organized into nine categories.
Each attack method opens a raw TCP socket to the target QEMU UART
port, sends a crafted payload, and parses the firmware response
to determine exploit success.
Key attack implementations include:
\begin{itemize}[leftmargin=*]
  \item \textbf{Buffer overflow} (3~variants) --- sends oversized
    \texttt{FW\_UPDATE} payloads (256--1\,024~bytes) targeting the
    unchecked 128-byte \texttt{fw\_update\_buf}; a heap-spray variant
    writes NOP sleds followed by a shellcode marker.
  \item \textbf{Authentication bypass} (2~variants) --- issues
    privileged commands (\texttt{FW\_UPDATE}, \texttt{REBOOT}) without
    a prior \texttt{AUTH} handshake; tests empty- and null-token
    acceptance.
  \item \textbf{Command injection} (2~variants) --- embeds shell
    metacharacters (\texttt{;}, \texttt{|}, backticks) and format
    strings (\texttt{\%x}, \texttt{\%n}) within \texttt{SET\_ID}
    payloads.
  \item \textbf{Firmware update exploitation} (2~variants) ---
    pushes unsigned binary payloads via \texttt{FW\_UPDATE} and
    triggers \texttt{FW\_APPLY} without integrity verification.
  \item \textbf{Information disclosure} (2~variants) --- reads
    auth tokens and firmware buffers via \texttt{DEBUG\_DUMP}; probes
    for verbose error messages leaking internal state.
  \item \textbf{Denial of service} (2~variants) --- rapid-fire
    command flooding (100~commands in burst) and resource exhaustion
    through repeated \texttt{FW\_UPDATE} writes.
  \item \textbf{State manipulation}, \textbf{replay}, and
    \textbf{protocol fuzzing} round out the suite with
    unauthorized state changes, captured-command replay, and
    random/boundary-value input generation.
\end{itemize}
Each attack returns an \texttt{AttackResult} dataclass containing
the exploit outcome (success/failure), CVSS~v3.1 severity level,
captured evidence, wall-clock duration, impact description, and
suggested mitigation.

\paragraph{Network Attacks (Suite~2).}
The \texttt{NetworkAttackFramework} orchestrates 14~network-level
attacks across five attack suites:
\begin{itemize}[leftmargin=*]
  \item \textbf{MQTT suite} (3) --- unauthorized topic subscription,
    malicious message injection to device command topics, and
    wildcard topic hijacking.
  \item \textbf{TCP/IP suite} (3) --- TCP connection hijacking via
    sequence-number prediction, man-in-the-middle proxy insertion,
    and SYN flood denial of service.
  \item \textbf{Protocol suite} (3) --- simulated Zigbee frame
    replay, simulated network-key extraction (demonstrating the
    well-known Trust Center link key weakness), and protocol
    downgrade verification (testing whether devices accept
    unencrypted connections).
  \item \textbf{Infrastructure suite} (3) --- ARP cache poisoning
    (conceptual on Docker bridge), simulated DNS response spoofing,
    and simulated WiFi deauthentication (modeled without
    physical-layer frame injection).
  \item \textbf{Traffic analysis suite} (2) --- simulated evil twin
    access point with credential harvesting, and active traffic
    fingerprinting via device identification commands.
\end{itemize}

\paragraph{Phantom-Delay Attack (Suite~3).}
Three attack variants adapted from Fu~et~al.~\cite{phantom-delay}
target IoT timeout behaviors via transparent TCP proxying:
state-update delay, erroneous automation execution, and routine
invalidation.  Each variant intercepts UART-over-TCP and MQTT
traffic on the Docker bridge network to selectively delay events
without triggering device-offline alerts.

\paragraph{Malicious SmartApp (Suite~4).}
A standalone attack module (\texttt{smartapp.py}, $\sim$380 lines)
targets the SmartThings Schema Connector's HTTP API.  The attack
uses raw HTTP (no external dependencies) to:
\begin{enumerate}[leftmargin=*]
  \item Probe the connector's \texttt{/health} endpoint.
  \item Send a \texttt{discoveryRequest} to enumerate all registered
    devices without valid OAuth credentials.
  \item Issue a \texttt{stateRefreshRequest} to steal the callback
    token stored in memory.
  \item Inject a forged \texttt{commandRequest} (e.g., DISARM) that
    the connector accepts without user confirmation.
\end{enumerate}
This attack targets CWE-269 (Improper Privilege Management) and
CWE-862 (Missing Authorization), with a CVSS base score of~8.8.

\paragraph{ESP32 Buffer Overflow (Suite~5).}
A standalone attack module (\texttt{esp32\_overflow.py},
$\sim$400~lines) targets the 128-byte command buffer in the ESP32
firmware.  The exploit constructs a 137-byte payload:
\begin{itemize}[leftmargin=*]
  \item 32~bytes of Thumb-mode NOP sled (\texttt{0xBF00} repeated).
  \item A 12-byte \texttt{VESPER\_PWNED} shellcode marker.
  \item 88~bytes of padding to reach the saved link register at
    offset~132 from the buffer start.
  \item A 4-byte fake return address (\texttt{0x20000081}, with
    Thumb bit set) pointing into the NOP sled in SRAM.
\end{itemize}
The payload triggers a device crash (DoS) when the overwritten
return address causes an invalid instruction fetch.  The device
becomes unresponsive for $\sim$12\,s before rebooting.  This attack
targets CWE-120 (Buffer Copy without Size Check) and CWE-787
(Out-of-bounds Write), with a CVSS base score of~9.8 (Critical).

\subsection{Security Evaluation Pipeline}

A \texttt{SecurityEvaluator} module (900~lines) processes raw attack
results into conference-grade metrics.  It computes CVSS~v3.1 base
scores using the standard formula (exploitability $\times$ impact),
maps each finding to MITRE ATT\&CK for IoT techniques (15~techniques
across 10~of 12~tactics), and classifies attacks along the seven-stage
IoT Cyber Kill Chain.  Statistical significance is assessed with
chi-squared, Fisher's exact, Kruskal-Wallis, and Mann-Whitney~U
tests.  The pipeline auto-generates six \LaTeX{} tables and six
publication-ready PDF figures (heatmaps, box plots, bar charts)
for direct inclusion in the manuscript.

\subsection{Configuration and Reproducibility}

All simulation parameters are specified in a single YAML
configuration file loaded by a Pydantic configuration model.
Key parameters include the random seed (\texttt{seed: 42}),
simulation duration, device inventory, persona selection, and LLM
model.  Experiment configurations for all evaluations are provided
in the repository under \texttt{vesper/evaluation/configs/}.
